\section{Corpus Map Across 82 Study Families}\label{sec:corpus-map}
\subsection{Origins, receiving projections and threat polarities}
Across all \FinalCore{} core families in our unified repository, 38 inherit runtime, serving or hardware channels (Origin~B), while 20 concern agent-workflow channels arising from orchestration, tool dispatch, browser navigation and multi-agent protocols (Origin~A; Table~\ref{tab:core-family-map}). Signal-generation origin is distinct from receiver location and threat polarity. The full map makes these distinctions visible family by family.
\begin{table*}[t]\centering\footnotesize
\caption{Complete map of the 58 core study families. Each family appears once; multiple layer/type labels describe different components or objectives within the family.}\label{tab:core-family-map}
\setlength{\tabcolsep}{3pt}
\begin{tabularx}{\textwidth}{@{}Ycc@{\hspace{1.5em}}Ycc@{}}\toprule
Core family & $G\,\to\,O$ & Type & Core family & $G\,\to\,O$ & Type\\\midrule
ArchSnoop~\cite{weng2026archsnoop} & $1\,\to\,1$ & \texttt{II} & PromptPeek~\cite{promptpeek2025} & $2\,\to\,4$ & \texttt{I}\\
Behind Bars~\cite{gu2026behindbars} & $1\,\to\,1$ & \texttt{I} & Remote keylogging~\cite{weiss2024prompt} & $2\,\to\,3$ & \texttt{I}\\
Edge EM leakage~\cite{alhalawani2026physicaledge} & $1\,\to\,1$ & \texttt{I} & Remote timing~\cite{carlini2024remote} & $2\,\to\,3$ & \texttt{I}\\
Energon~\cite{chaudhuri2025energon} & $1\,\to\,1$ & \texttt{II} & Routing metadata~\cite{chen2026routingmetadata} & $2\,\to\,4$ & \texttt{I}\\
I Know What You Said~\cite{gao2025iknow} & $1\,\to\,1$ & \texttt{I} & SafeKV~\cite{chu2025safekv} & $2\,\to\,4$ & \texttt{I}\\
MoEcho~\cite{ding2025moecho} & $1\,\to\,1$ & \texttt{I} & Speculative decoding~\cite{wei2024speculation} & $2\,\to\,3$ & \texttt{I}\\
PAGE404~\cite{nazaraliyev2026page404} & $1\,\to\,1$ & \texttt{I} & SpliceLeak~\cite{spliceleak2026} & $2\,\to\,4$ & \texttt{I}\\
Spill the Beans~\cite{spillbeans2025} & $1\,\to\,1$ & \texttt{I} & SplitDecoy~\cite{politis2026splitdecoy} & $2\,\to\,2$ & \texttt{I}\\
Detokenization Leaks~\cite{weiss2026detokenization} & $1,2\,\to\,1$ & \texttt{I} & Time Will Tell~\cite{timewilltell2024} & $2\,\to\,4$ & \texttt{I}\\
LLM-serving RCA~\cite{scheinert2026rca} & $1,2\,\to\,2$ & \texttt{IV} & Whisper Leak~\cite{whisperleak2025} & $2\,\to\,3$ & \texttt{I}\\
LLM-WFIN~\cite{jiao2025llmwfin} & $1,2\,\to\,1,3$ & \texttt{I} & MIRAGE~\cite{revankar2026mirage} & $2,5\,\to\,2$ & \texttt{IV}\\
OTRO~\cite{lee2026otro} & $1,2\,\to\,1$ & \texttt{I} & Cloud traffic benchmark~\cite{pouryousef2026traffic} & $2,5,8\,\to\,3$ & \texttt{I,II}\\
Shape and Substance~\cite{hadad2026shape} & $1,2\,\to\,1$ & \texttt{I} & Web-agent unmasking~\cite{fayolle2026webagents} & $3,7\,\to\,3,7$ & \texttt{II,III}\\
SparLeak~\cite{chen2026sparleak} & $1,2\,\to\,1$ & \texttt{I} & LumiMAS~\cite{solomon2026lumimas} & $4,5\,\to\,4,6$ & \texttt{IV}\\
SparSEEty~\cite{jo2026sparseety} & $1,2\,\to\,1$ & \texttt{I} & AgentPrint~\cite{agentprint2025} & $5\,\to\,3$ & \texttt{I,II}\\
TDXRay~\cite{hornetz2026tdxray} & $1,2\,\to\,1$ & \texttt{I} & Cycle detection~\cite{george2025cycles} & $5\,\to\,4,6$ & \texttt{IV}\\
ClawGuard~\cite{gan2026clawguard} & $1,5\,\to\,1$ & \texttt{IV} & AgentProv~\cite{wang2026agentprov} & $5,6\,\to\,6$ & \texttt{II}\\
Cache contention~\cite{abubakar2026contention} & $2\,\to\,4$ & \texttt{I} & Lidar~\cite{wang2026lidar} & $5,6\,\to\,6$ & \texttt{II}\\
CachePrune~\cite{wu2026cacheprune} & $2\,\to\,4$ & \texttt{I} & ProcGrep~\cite{oderinwale2026procgrep} & $5,6\,\to\,4,6$ & \texttt{II}\\
Early Bird~\cite{earlybird2025} & $2\,\to\,4$ & \texttt{I} & SIGLEAK~\cite{geng2026sigleak} & $5,6\,\to\,4,6$ & \texttt{I}\\
InputSnatch~\cite{inputsnatch2024} & $2\,\to\,4$ & \texttt{I} & Campaign attribution~\cite{park2026crosscampaign} & $5,6,8\,\to\,6,8$ & \texttt{II}\\
KVGov~\cite{addagada2026kvgov} & $2\,\to\,2,4$ & \texttt{I} & Dangol workflow~\cite{dangol2026workflow} & $5,6,8\,\to\,3,6,8$ & \texttt{I}\\
Leaky Language Models~\cite{majidi2026leakylms} & $2\,\to\,4$ & \texttt{I} & AgentTell~\cite{shahriar2026agenttell} & $5,7\,\to\,7$ & \texttt{I}\\
LLMs Have Rhythm~\cite{alhazbi2025rhythm} & $2\,\to\,3$ & \texttt{II} & AHB~\cite{zhu2026turing} & $5,7\,\to\,7$ & \texttt{III}\\
ML privacy channels~\cite{debenedetti2024privacy} & $2\,\to\,4$ & \texttt{I} & FP-Agent~\cite{wang2026fpagent} & $5,7\,\to\,7$ & \texttt{II}\\
MoE prompt stealing~\cite{yona2024stealing} & $2\,\to\,4$ & \texttt{I} & PersonaFingerprint~\cite{song2026personafingerprint} & $5,7\,\to\,3$ & \texttt{I}\\
NetEcho~\cite{netecho2025} & $2\,\to\,3$ & \texttt{I} & UI Traces~\cite{uitraces2026} & $5,7\,\to\,7$ & \texttt{II}\\
Prompt Echo~\cite{prompt_echo2026} & $2\,\to\,3$ & \texttt{I} & Whose Agent?~\cite{kang2026whoseagent} & $5,7\,\to\,7$ & \texttt{II}\\
Prompt-cache audit~\cite{gu2025promptcaching} & $2\,\to\,4$ & \texttt{I} & WRA~\cite{jeong2026wra} & $5,7\,\to\,3$ & \texttt{I}\\
\bottomrule\end{tabularx}
\par\smallskip\raggedright $G$ denotes signal-generation layers and $O$ the receiving observation layers. Layer indices: 1 hardware; 2 serving/runtime; 3 transport; 4 application/API; 5 orchestration; 6 tools; 7 browser/environment; 8 multi-agent communication. Types: I confidentiality/privacy; II attribution; III human--agent distinction/evasion; IV authorized diagnostics. These are objectives and observer roles, not mutually exclusive security guarantees. Privileged operator receivers and physical receivers are marked by the layer they inspect. Source versions, locators and the normalization of prose-only layer descriptions are in the artifact.
\end{table*}


\paragraph{Cross-origin synthesis}
Origin B includes token/packet lengths, cache timing and co-resident hardware channels; Origin A includes visited domains, tool dispatch and browser actions. Receiver permissions cut across these origins: a remote network observer, a same-site script and a privileged diagnostic hook expose different views even when they share an attribution target. Thus Origin A/B cannot substitute for Types I--IV or observer permissions. WRA's read-only privacy decoys, AHB's stateful evasion and SplitDecoy's exposed return path illustrate transfers that fail when these axes are collapsed \cite{jeong2026wra,zhu2026turing,politis2026splitdecoy}. The map retains every core family rather than using these examples as a substitute for coverage.

Layer incidence highlights where the systematization's comparison problem arises. In the 58-family map, signal-generation labels cluster at P2 serving/runtime (32 families) and P5 orchestration (19); receiving views cluster at P4 application/API (18), P1 hardware (16), and P3 transport (14). Generation at P3/P4 is rare (one family each), despite frequent observation at those layers. This source/receiver asymmetry makes the transfer question concrete: protection at a frequent source layer cannot be presumed to cover a different receiving projection. The map records layer sets, so it does not assign a causal edge to each source/receiver pair.

Purpose labels further show why a single “agent side channel” category is insufficient: Type I appears 41 times, Type II 13, Type III 2, and Type IV 5. The 32 map-only families extend coverage beyond mitigation comparisons: 15 list serving origins, 11 hardware and 11 orchestration; 19 concern confidentiality, nine attribution, one evasion and all five authorized-diagnostic families. They comprise 22 families without a scoped mitigation evaluation, five diagnostics, four unbound quantitative assertions, and SplitDecoy's boundary case (\S\ref{sec:defenses}). Mapping these families locates privacy targets, operator capabilities and candidate validation environments even when no scoped mitigation record enters the ledger. Counts are overlapping family-label incidences; the audit then applies its reuse rules to the 39 scoped result records.

\subsection{Relationship to prior systematizations}
Prior surveys establish compound-system threat models and confidential-computing boundaries. Banerjee et al.\ categorize cross-layer compound AI threats under MITRE tactics; Forough et al.\ distinguish TEE isolation boundaries from side-channel resistance and note that inference protections need not cover agent memory or action coordination \cite{banerjee2024compound,forough2026secrets}. Kotevska et al.\ propose a six-element privacy contract (asset, observer, channel, permitted disclosure, guarantee, evidence), with composition conditions framed as formal conjectures \cite{kotevska2026privacyfoundations}. Classical website-fingerprinting studies establish bandwidth/latency costs for transport shaping, including Tamaraw, WTF-PAD, FRONT and Walkie-Talkie \cite{cai2014systematic,juarez2016wtfpad,gong2020front,wang2017walkietalkie}. Our contribution is to make reuse decisions inspectable across 82 mapped families and 39 source-bound result records, then organize the empirical checks around four boundaries: emitter and transport; stateful execution and the exposed link; repeated real-agent trajectories; and live timing with adapted website-fingerprinting controls. This evidence map shows why a defense's name or layer is insufficient for transfer: the receiver projection, state effects, workload and cost object must also align. Our experiments test these specific comparison obligations on a loopback TLS path and a pinned stateful tool simulator. These measurements connect established mechanisms to the cost and residual inference observed in streamed LLM and agent traffic.
